\section{Averaging as an Approximate Weighted LMO}
\label{app:matnormal-01}

The single-batch LMO in \Cref{sec:sample-loss-lmo} has two inputs, the
gradient in its objective and the per-sample second moment in its
outergradient constraint. To draw on recent batches while keeping the LMO
local to the current iterate, we evaluate every recent sample's gradient at
$W_t$ and average the two inputs across those samples, giving the ideal
moments $\overline G_t$ and $\overline\Sigma_t$. This appendix explains how
\Cref{alg:ours}
approximates them from the LoRA factor gradients already computed by
backpropagation.

\Cref{sec:weighted-lmo-01} writes the weighted LMO and isolates the two
moments through which past steps enter the problem.
\Cref{sec:weighted-estimation-01}
replaces those moments by quantities built from $G_{A,s}$ and $G_{B,s}$,
giving the $p,q$ recursion in \Cref{alg:ours} and the
update~\eqref{eq:polora-update}.
\Cref{sec:scale-placement-01} explains why \Cref{alg:ours} normalizes the
accumulated vectors after averaging rather than normalizing each estimate
before it enters the average. The update is invariant to the scale of the curvature factors,
but the average that estimates them depends on that scale.

\subsection{The Weighted LMO}
\label{sec:weighted-lmo-01}

Let $s$ index optimizer steps, and let $i$ index samples in the
minibatch of size $n_s$ at step $s$. Write
$G_{s,i}:=\nabla\ell_{x_{s,i}}(W_s)$,
$g_{s,i}:=\vec(G_{s,i})$, and
$G_s:=n_s^{-1}\sum_iG_{s,i}$. Let $m_{t,s}$ and $w_{t,s}$ be
nonnegative weights on past steps, one for the objective and one for
the constraints.

In the ideal weighted LMO, every sample from a step $s\le t$ contributes
the gradient it would have at the current weights $W_t$, so drawing on past
steps changes only which samples enter the averages, not where their
gradients are evaluated. The two resulting moments are
\[
\overline G_t
  =
  \sum_{s\le t}\frac{m_{t,s}}{n_s}\sum_{i=1}^{n_s}
  \nabla\ell_{x_{s,i}}(W_t),
\qquad
\overline\Sigma_t
  =
  \sum_{s\le t}\frac{w_{t,s}}{n_s}\sum_{i=1}^{n_s}
  \vec(\nabla\ell_{x_{s,i}}(W_t))
  \vec(\nabla\ell_{x_{s,i}}(W_t))^\top .
\]
The first moment is the gradient in the linear objective. For the
constraint, the relevant gradients are rescaled as
$\sqrt{w_{t,s}/n_s}\,\nabla\ell_{x_{s,i}}(W_t)$, and their outer
products sum to $\overline\Sigma_t$. Applying the leverage argument of
\Cref{lem:gradprop} to these rescaled gradients gives a relaxation of
the outergradient constraint with the same form
as~\eqref{eq:polorafull-hard}. The relaxed problem is
\begin{equation}\label{eq:weighted-relaxed-01}
 \begin{array}{ll}
  \underset{\Delta W}{\mathrm{minimize}} &
  \ip{\overline G_t, \Delta W} \\[2pt]
  \mbox{subject to} &
  \displaystyle
  \max_{\widetilde{G}\in\mathcal{G}(\overline\Sigma_t)}
  \ip{\widetilde{G},\Delta W}\le\tau .
  \end{array}
\end{equation}
\Cref{alg:ours} stores neither directly, because it never recomputes old
samples' gradients at $W_t$ and the per-sample second moment is too large
to keep.

\paragraph{The exponential weights.}
If every past gradient could be recomputed at $W_t$, equal weights would be
the direct extension of the single-batch LMO to samples from earlier steps.
Since that average is unavailable, \Cref{alg:ours} uses two geometrically decaying windows,
one for the objective and one for the constraint. These windows use
constant memory and put less weight on older gradients and factors, which
were computed farther from the current iterate.

\Cref{alg:ours} uses look-ahead momentum for the first moment,
\[
  M_{t+1}=\beta_1M_t+(1-\beta_1)G_t,
  \qquad
  \widehat M_t=\beta_1M_{t+1}+(1-\beta_1)G_t,
\]
and an EMA for the second moment. Unrolling these recursions
gives the induced step weights
\begin{equation}\label{eq:EMAweights-01}
  m_{t,s}=\beta_1^2(1-\beta_1)\beta_1^{t-1-s}
  \quad(s<t),
  \qquad
  m_{t,t}=1-\beta_1^2,
  \qquad
  w_{t,s}=(1-\beta_2)\beta_2^{t-s}.
\end{equation}
We write $\E_{s\le t}[\phi(G_s)]=\sum_{s\le t}w_{t,s}\,\phi(G_s)$ for an
average over past steps with weights $w_{t,s}$. Replacing each unavailable
current-weight gradient in $\overline G_t$ by the observed batch gradient
$G_s$ turns the first moment into the look-ahead momentum $\widehat M_t$,
with the objective weights $m_{t,s}$ in~\eqref{eq:EMAweights-01}. The
optimizer never forms the $d_{\mathrm{out}}\times d_{\mathrm{in}}$ matrix
$\widehat M_t$, so the next subsection expresses its needed projections and
the curvature EMA in terms of factor gradients.

\subsection{Estimating the Moments from Factor Gradients}
\label{sec:weighted-estimation-01}

The optimizer maintains momenta for $G_{A,s},G_{B,s}$ and diagonal
curvature vectors instead of the dense moments $\overline G_t$ and
$\overline\Sigma_t$.

\paragraph{From moments at $W_t$ to momentum and curvature variables.}
The replacement from the ideal moments to the quantities in
\Cref{alg:ours} uses four approximations.
\begin{enumerate}[leftmargin=1.6em,itemsep=2pt,topsep=2pt]
  \item A gradient of an old sample at $W_t$ is replaced by the gradient
  observed when that sample was processed,
  $\nabla\ell_{x_{s,i}}(W_t)\leadsto\nabla\ell_{x_{s,i}}(W_s)$.
  \item The second moment within the batch is approximated by the outer
  product of the batch gradient,
  \[
    \frac1{n_s}\sum_i g_{s,i}g_{s,i}^\top
    \quad\leadsto\quad
    g_sg_s^\top,
    \qquad g_s=\vec(G_s).
  \]
  For $n_s>1$, this replaces the per-sample second moment by the
  second moment of the minibatch gradient, as in AdaGrad~\cite{adagrad},
  Adam~\cite{adam}, and Shampoo~\cite{shampoo}.
  \item The resulting dense second moment is modeled by diagonal
  Kronecker factors $P^\star_t,Q^\star_t$ that the running estimates $P,Q$
  of \Cref{alg:ours} approximate,
  \begin{equation}\label{e-diag-kron-second-mom-model-01}
    \E_{s\le t}[g_sg_s^\top]\approx\kappa_t\,(Q^\star_t\otimes P^\star_t).
  \end{equation}
  We fix $\norm{\diag P^\star_t}_\infty=\norm{\diag Q^\star_t}_\infty=1$ and
  let the scalar $\kappa_t>0$ carry the magnitude, which does not affect the
  final normalized direction.
  \item Moments of $G_{A,s}$ and $G_{B,s}$ use the current LoRA factors inside the
  recency-weighted sums,
  $A_s\leadsto A_t$ and $B_s\leadsto B_t$ for $s\le t$. The exponential
  weights make this a local approximation by downweighting older steps.
\end{enumerate}
The first three approximations reduce the constraint moment to
\[
  \overline\Sigma_t\approx\E_{s\le t}[g_sg_s^\top],
\]
which the diagonal Kronecker model~\eqref{e-diag-kron-second-mom-model-01}
writes as $\kappa_t(Q^\star_t\otimes P^\star_t)$.

\paragraph{Factor momenta.}
LoRA backpropagation already computes the factor gradients
$G_{A,s}=B_s^\top G_s$ and $G_{B,s}=G_sA_s^\top$. Forming the dense
gradient $G_s$ would require a $d_{\mathrm{out}}\times d_{\mathrm{in}}$
matrix per layer. The optimizer therefore represents the first moment only
through the two projections needed by the factor update. Under the
linearized product update~\eqref{eq:DeltaW-lin}, the objective of
\eqref{eq:weighted-relaxed-01} decomposes as
\[
  \ip{\widehat M_t,\,\Delta B\,A+B\,\Delta A}
  =
  \ip{\widehat M_tA^\top,\Delta B}
  +
  \ip{B^\top\widehat M_t,\Delta A}.
\]
The direction step then needs only $\widehat M_tA_t^\top$ and
$B_t^\top\widehat M_t$. The factor momenta in \Cref{alg:ours} average
$G_sA_s^\top$ and $B_s^\top G_s$. After replacing $A_s,B_s$ by
$A_t,B_t$ inside the weighted sums, these approximate the desired
projections $\widehat M_A\approx B_t^\top\widehat M_t$ and
$\widehat M_B\approx \widehat M_tA_t^\top$, up to the drift of the factors
across the EMA window.

\paragraph{Moments of the factor gradients.}
With $\E_{s\le t}[g_sg_s^\top]$ as the dense second moment, replacing
$A_s,B_s$ by $A_t,B_t$ links the second moments of the factor gradients to
the weighted second moment of $G_s$. For the $A$ factor,
$\vec(G_{A,s})=(I\otimes B_s^\top)\vec(G_s)$, so
\begin{align*}
 \Sigma_{A,t}
 &:= \E_{s\le t}\big[\vec(G_{A,s})\vec(G_{A,s})^\top\big] \\
 &= \E_{s\le t}\big[
      (I\otimes B_s^\top)g_sg_s^\top(I\otimes B_s)
    \big] \\
 &\approx
    (I\otimes B_t^\top)
    \E_{s\le t}\big[g_sg_s^\top\big]
    (I\otimes B_t) \\
 &\approx
    \kappa_t\,(I\otimes B_t^\top)
    (Q^\star_t\otimes P^\star_t)
    (I\otimes B_t) \\
 &= \kappa_t\,Q^\star_t\otimes(B_t^\top P^\star_tB_t).
\end{align*}
The first approximation replaces $B_s$ by $B_t$ inside the weighted
sum, the second applies the diagonal Kronecker model, and the final
equality is the mixed product rule. The analogous calculation for
$G_{B,s}$, written with $\vec(G_{B,s}^\top)$ so that the Kronecker
blocks are indexed by the output coordinate, gives
\begin{align}
  \Sigma_{A,t}
    &\approx \kappa_t\,Q^\star_t\otimes C^\star_{A,t},
  \nonumber \\
  \Sigma_{B,t}
    &:= \E_{s\le t}\big[
      \vec(G_{B,s}^\top)\vec(G_{B,s}^\top)^\top
    \big]
    \approx \kappa_t\,P^\star_t\otimes C^\star_{B,t},
  \label{eq:observable-moments-01}
\end{align}
where
$C^\star_{A,t}=B_t^\top P^\star_tB_t$ and
$C^\star_{B,t}=A_tQ^\star_tA_t^\top$.

\paragraph{Fitting the diagonal factors.}
We fit the diagonal factors to the factor-gradient moments in
\eqref{eq:observable-moments-01}. In this paragraph, all quantities are
at the current step, so we suppress the subscript $t$. Following
KL-Shampoo~\cite{klshampoo}, the fit uses the log-determinant
divergence~\cite{logdetdiv}
\[
  D(\Sigma,S)=\Tr(\Sigma S^{-1})-\log\det(\Sigma S^{-1})-d,
\]
defined for positive definite matrices of common dimension $d$.
The $A$-gradient moment identifies $Q^\star$ once
$C^\star_A=B^\top P^\star B$ is fixed, and the $B$-gradient moment
identifies $P^\star$ once $C^\star_B=AQ^\star A^\top$ is fixed. Each fit
therefore freezes the opposite factor at its current estimate:
\begin{equation}\label{eq:factor-fit-01}
  \hat{q}(\widehat P)
  =
  \argmin_{q>0}
  D\bigl(\Sigma_A,\diag(q)\otimes\widehat C_A\bigr),
  \qquad
  \hat{p}(\widehat Q)
  =
  \argmin_{p>0}
  D\bigl(\Sigma_B,\diag(p)\otimes\widehat C_B\bigr),
\end{equation}
where
$\widehat C_A=B^\top\widehat P B$ and
$\widehat C_B=A\widehat Q A^\top$ for diagonal
$\widehat P,\widehat Q\succ0$. The blockwise first-order conditions give
the closed forms
\[
  \hat{q}(\widehat P)
  =
  \frac1r\diag\!\left(
    \E_{s\le t}[G_{A,s}^\top\widehat C_A^{-1}G_{A,s}]
  \right),
  \qquad
  \hat{p}(\widehat Q)
  =
  \frac1r\diag\!\left(
    \E_{s\le t}[G_{B,s}\widehat C_B^{-1}G_{B,s}^\top]
  \right).
\]

For fixed positive definite $\Sigma$ and $S$, define the best scalar
multiple
\[
  c^\star(\Sigma,S)
  :=
  \argmin_{c>0}D(\Sigma,cS)
  =
  \Tr(\Sigma S^{-1})/d .
\]

\begin{lemma}[Consistency of the fit from factor gradients]\label{lem:consistency-01}
Assume the moments of the factor gradients satisfy
\[
  \Sigma_A=\kappa\,Q^\star\otimes C^\star_A,
  \qquad
  \Sigma_B=\kappa\,P^\star\otimes C^\star_B
\]
exactly, with
$C^\star_A=B^\top P^\star B$,
$C^\star_B=AQ^\star A^\top$, and
$\norm{\diag P^\star}_\infty=\norm{\diag Q^\star}_\infty=1$.
For any diagonal $\widehat P\succ0$ and $\widehat Q\succ0$, the
fits~\eqref{eq:factor-fit-01} return
\[
  \hat{q}(\widehat P)
  =
  \kappa\,c^\star(C^\star_A,\widehat C_A)\,\diag(Q^\star),
  \qquad
  \hat{p}(\widehat Q)
  =
  \kappa\,c^\star(C^\star_B,\widehat C_B)\,\diag(P^\star).
\]
Therefore
$\hat{q}/\norm{\hat{q}}_\infty=\diag(Q^\star)$ and
$\hat{p}/\norm{\hat{p}}_\infty=\diag(P^\star)$.
\end{lemma}
\begin{proof}
Under the stated moment model, $\Sigma_A$ is block diagonal
with $j$th block $\kappa q^\star_j C^\star_A$, while
$\diag(q)\otimes\widehat C_A$ has $j$th block $q_j\widehat C_A$.
Traces and log-determinants add across blocks, so the objective for
$q$ is
\[
  \sum_j D(\kappa q^\star_j C^\star_A, q_j\widehat C_A).
\]
The $j$th term is minimized by its best scalar multiple,
\[
  \hat q_j
  =
  c^\star(\kappa q^\star_j C^\star_A,\widehat C_A)
  =
  \kappa q^\star_j c^\star(C^\star_A,\widehat C_A).
\]
The scalar $\kappa c^\star(C^\star_A,\widehat C_A)$ is independent of
$j$, and the normalization removes it. The proof for $\hat p$ is the
same with the two factors exchanged.
\end{proof}

\paragraph{Online estimator.}
The weighted fit in~\eqref{eq:factor-fit-01} would combine all past
factor gradients using the same current matrices. The online recurrence
instead evaluates each term once, using the matrices
$C_{A,t}=B_t^\top P_tB_t$ and
$C_{B,t}=A_tQ_tA_t^\top$ available at that step:
\begin{equation}\label{eq:ours-klcoupling-01}
\begin{aligned}
  \hat{q}_t(P_t)
    &= \frac1r\diag\!\left(G_{A,t}^\top C_{A,t}^{-1}G_{A,t}\right),
  &
  \hat{p}_t(Q_t)
    &= \frac1r\diag\!\left(G_{B,t}C_{B,t}^{-1}G_{B,t}^\top\right),
  \\
  q_{t+1}
    &= \beta_2q_t+(1-\beta_2)\hat{q}_t(P_t),
  &
  p_{t+1}
    &= \beta_2p_t+(1-\beta_2)\hat{p}_t(Q_t),
  \\
  Q_{t+1}
    &= \diag\!\bigl(q_{t+1}/\norm{q_{t+1}}_\infty\bigr),
  &
  P_{t+1}
    &= \diag\!\bigl(p_{t+1}/\norm{p_{t+1}}_\infty\bigr).
\end{aligned}
\end{equation}
Unrolling the EMA, up to the initialization tail, gives
\[
\begin{aligned}
  q_{t+1}
  &=
  \sum_{s\le t}w_{t,s}\hat q_s(P_s)
  \approx
  \frac1r\diag\!\left(
    \sum_{s\le t}w_{t,s}\,G_{A,s}^\top C_{A,t}^{-1}G_{A,s}
  \right)
  =
  \hat q(P_t),\\
  p_{t+1}
  &=
  \sum_{s\le t}w_{t,s}\hat p_s(Q_s)
  \approx
  \frac1r\diag\!\left(
    \sum_{s\le t}w_{t,s}\,G_{B,s}C_{B,t}^{-1}G_{B,s}^\top
  \right)
  =
  \hat p(Q_t).
\end{aligned}
\]
The exact recurrence is the EMA over the per-step fits
$\hat q_s(P_s),\hat p_s(Q_s)$. Reading this EMA as the weighted fit at
the current matrices replaces the matrices inside older terms by
$C_{A,t}$ and $C_{B,t}$, using the same slow-evolution approximation as
above. Under the moment model,
\Cref{lem:consistency-01} then predicts
\[
  Q_{t+1}
  =
  \diag\!\bigl(q_{t+1}/\norm{q_{t+1}}_\infty\bigr)
  \approx Q^\star_t,
  \qquad
  P_{t+1}
  =
  \diag\!\bigl(p_{t+1}/\norm{p_{t+1}}_\infty\bigr)
  \approx P^\star_t .
\]

\paragraph{Relation to KL-Shampoo.}
KL-Shampoo~\cite{klshampoo} also uses the log-determinant divergence to
fit Kronecker factors. The fit here differs in three ways.
\begin{itemize}[leftmargin=1.5em,itemsep=1pt,topsep=2pt]
  \item \emph{No distributional assumptions.} KL-Shampoo fits the second
  moment of a stationary gradient distribution. Here $\E_{s\le t}$ is a
  weighted sum over the realized gradients, not an expectation. The
  consistency claim is conditional on the displayed moment equations, not
  on a distributional limit.
  \item \emph{The fitted moments are factor-gradient moments.} KL-Shampoo
  fits Kronecker factors to the second moment of the weight gradient $G$.
  Here $G$ is never formed. Each diagonal factor is fit from the second
  moment of its corresponding LoRA factor gradient, and the two fits are
  coupled through $C_A$ and $C_B$.
  \item \emph{Fixed-moment statement.} KL-Shampoo
  justifies its moving average as a stochastic proximal gradient step on the
  divergence, correct in the limit. Here \Cref{lem:consistency-01} is an
  algebraic statement for fixed moments. Conditional on the moment model,
  using an inexact estimate of the other factor only rescales the fitted
  vector, which the normalization removes, and the EMA is exact over the
  per-step fits.
\end{itemize}

\paragraph{Recovering the update of \Cref{alg:ours}.}
Substituting the fitted model $\kappa_t(Q_t\otimes P_t)$ for
$\overline\Sigma_t$ and the look-ahead momentum approximation for
$\overline G_t$ yields the preconditioned spectral
LMO~\eqref{eq:polorafull} (\Cref{lem:pseudogradient-spectral}). The
direction and magnitude steps of \Cref{sec:sample-loss-lmo}, with the
factor momenta standing in for the two dense momentum projections, give
the update~\eqref{eq:polora-update}. \Cref{alg:ours} updates the
curvature factors after the direction step, so step $t$ uses the estimate
from the previous step. This avoids forming and inverting $C_A$ and
$C_B$ twice per optimizer step.

\subsection{Normalization Placement}
\label{sec:scale-placement-01}

\methodname{} normalizes the accumulated vectors $q_{t+1},p_{t+1}$ when
forming $Q_{t+1},P_{t+1}$ in~\eqref{eq:ours-klcoupling-01}. It does not
normalize each estimate $\hat{q}_t,\hat{p}_t$ before adding it to
the EMA. The update is invariant to the overall scale of $P$ and $Q$, but
the EMA that estimates them is not.

\paragraph{Rescaling does not change the update.}
The update depends on $P$ and $Q$ only through their direction. Under a
rescaling $Q\mapsto cQ$ with $c>0$, both $Q^{-1/2}$ and
$C_B^{-1/2}=(AQA^\top)^{-1/2}$ pick up a factor $c^{-1/2}$. Inside the
matrix sign this scalar has no effect, since the sign is unchanged by
positive rescaling. Outside it, the scalar multiplies $D_A$ and $D_B$ by
$c^{-1/2}$, which the division by $\norm{D_A}_2$ and
$\norm{D_B}_2$ in the magnitude step cancels. The factor $P$ enters in
the same way. The update therefore sees the direction of the curvature
factors, but not the magnitude $\kappa_t$ of the model
\eqref{e-diag-kron-second-mom-model-01}. Normalizing the accumulated
vectors leaves the update unchanged.

\paragraph{Rescaling changes the average.}
The coupled fit makes each estimate at step $t$ scale inversely with the
curvature factor used to compute it,
\[
  \hat{q}_t(aP_t)=a^{-1}\,\hat{q}_t(P_t),
  \qquad
  \hat{p}_t(aQ_t)=a^{-1}\,\hat{p}_t(Q_t)
  \qquad(a>0),
\]
since $C_A=B^\top PB$ is linear in $P$, and likewise $C_B$ in $Q$.
Normalizing after averaging gives every factor $P_t,Q_t$ the same scale,
$\norm{\diag P_t}_\infty=\norm{\diag Q_t}_\infty=1$, so the estimates
enter the EMA on a common scale. If the factor used at step
$s$ carried a scale $a_s$ that varied with $s$, the identity above would scale
$\hat{q}_s$ by $1/a_s$, so it would enter with effective weight
$w_{t,s}/a_s$ rather than the intended weight $w_{t,s}$ in
\Cref{sec:weighted-lmo-01}. Normalizing the estimates before
averaging would also change the estimate, since normalization does not
commute with the mean. The same average-then-normalize order appears in
normalized SGD~\cite{cutkosky_nigt}, where momentum precedes
normalization, and in Muon~\cite{denoise_ortho}, where momentum precedes
orthogonalization.
